\section{Regular Languages}
We begin by consider a class of languages known as the \textbf{regular
languages}. These languages can be defined by \textit{finite automaton},
\textit{regular expressions}, or a reduced form of recursive notation known as
\textit{linear grammars}.\\

Intuitively, regular languages comprises the finite languages and infinite
languages that involve modulo counting.

\subsection{Deterministic Finite Automata}
\begin{definition}[DFA]
  A \textit{DFA} $A = \left( Q,\Sigma,\delta,q_0,F \right) $, where:
  \begin{itemize}
    \item $Q$: A finite set of states
    \item $\Sigma$: A finite set of input symbols
    \item $\delta$: A transition function taking as argument a state and an
      input symbol and returns a state. $\delta: Q\times \Sigma\to Q$
    \item $q_0$: A start state, $q_0\in Q$
    \item $F$: A set of final or accepting states. $F\in \mathcal{P}\left( Q \right) $
  \end{itemize}
\end{definition}
\begin{remark}
  In CS3231, we consider only DFA with finite set of states.
\end{remark}
Suppose $a_1a_2\cdots a_{n}$ a sequence of input symbols. Starting at the start
state $q_0$, we consult the transition function $\delta$, to find the state of
the DFA $\delta\left( q_0,a_1 \right) =q_1$ after processing the first input
symbol $a_1$.
\begin{example}
  To generate the language $L=\{ \omega | \omega$ is of the form $x01y$
  for some strings $x$ and $y$ consisting of $0$'s and $1$'s only$\}$,
  consider the DFA $A=\left( \left\{ q_0,q_1,q_2 \right\} ,\left\{ 0,1
  \right\},\delta,q_0,\left\{ q_1 \right\}\right) $, where $\delta$ is defined:
  \[\delta\left( q_0,0 \right) =q_2,\quad \delta\left( q_0,1 \right) =q_0,\quad
  \delta\left( q_1,0 \right) =\delta\left( q_1,1 \right) =q_1,\quad
\delta\left( q_2,0 \right) =q_2,\quad\delta\left( q_2,1 \right) =q_1\] 
  We have alternative notations for description of DFAs, namely the transition diagram
  and the transition table.
  \begin{figure}[h!]
    \centering
    \raisebox{-0.5\height}{\includegraphics[width=0.5\textwidth]{dfa-diag}}
    \hspace*{.2in}
    \raisebox{-0.5\height}{\includegraphics[width=0.25\textwidth]{dfa-table}}
    %\raisebox{-0.5\height}{\begin{tabular}{ c c c }
    %  \toprule
    %   & $0$ & $1$ \\
    %   \midrule
    %  $q_0$ & $q_2$ & $q_0$ \\
    %  $q_1$ & $q_1$ & $q_1$ \\
    %  $q_2$ & $q_2$ & $q_1$ \\
    %  \bottomrule
    %\end{tabular}}
  \end{figure}
\end{example}
To extend the transition function to strings, we define an \textbf{extended
transition function} $\hat{\delta}:Q\times \Sigma^{*}\to Q$ to describe the
state (given a starting state), after processing a sequence of input symbols.
We define by induction on the length of the input string:
\begin{enumerate}
  \item $\hat{\delta}\left( q,\epsilon \right) =q$. Starting in state $q$, if
    no inputs are read, no change in state.
  \item Suppose $w$ a string of the form $xa$, where $x\in \Sigma^{*},a\in \Sigma$,
    \[\hat{\delta}\left( q,xa \right) =\delta\left( \hat{\delta}\left( q,x \right) ,a \right).\] 
    That is, we first process all but the last symbol of $w$, then its last symbol.
\end{enumerate}
\begin{definition}[Language of a DFA]
  The \textit{language} of a DFA $A$ is denoted $L(A)$, and is defined by
  \[L(A)=\left\{ \omega | \hat{\delta}\left( q_0,\omega \right) \in F \right\},\]
  the set of strings $\omega$ that the start state $q_0$ may take to arrive at
  one of the accepting states.
\end{definition}

\subsection{Nondeterministic Finite Automata}
We now consider an automata in which we can be in several state at once.
In an NFA, it is not necessary that all states have all inputs symbols
defining a transition. In those situations, the thread of the NFA's existence
corresponding to those states simply ``dies''.
\begin{figure}[h!]
  \centering
  \includegraphics[width=0.5\textwidth]{nfa-diag}
\end{figure}
\begin{definition}[NFA]
  An \textit{NFA} $A = \left( Q,\Sigma,\delta,q_0,F \right) $, where:
  \begin{itemize}
    \item $Q$: A finite set of states
    \item $\Sigma$: A finite set of input symbols
    \item $\delta$: A transition function taking as argument a state and an
      input symbol and returns \textbf{a set of states}. $\delta: Q\times \Sigma\to \mathcal{P}(Q)$
    \item $q_0$: A start state, $q_0\in Q$
    \item $F$: A set of final or accepting states. $F\in \mathcal{P}\left( Q \right) $
  \end{itemize}
\end{definition}
\begin{remark}
  The only difference between an NFA and a DFA is the return type of the
  transition function.
\end{remark}
Similarly, we define the \textbf{extended transition function}
$\hat{\delta}:Q\times \Sigma^{*}\to\mathcal{P}\left( Q \right) $, the set of
states that the NFA will be in if it starts in a given state $q$ and processes
the string $\omega$. We define by induction on the length of the input string:
\begin{enumerate}
  \item $\hat{\delta}\left( q,\epsilon \right) =\left\{ q \right\}$. If no
    inputs are read, we are only in the state we began in.
  \item Suppose $w$ a string of the form $xa$, where $x\in \Sigma^{*},a\in
    \Sigma$, and that $\hat{\delta}\left( q,x \right) =\left\{ p_1,p_2,\ldots,p_k \right\}$
    \[\hat{\delta}\left( q,w \right) =\bigcup_{i=1}^{k}\delta\left( p_i,a
    \right) =\left\{ r_1,r_2,\ldots,r_m \right\}\]
    That is, we first compute $\hat{\delta}\left( q,x \right)$, then follow any
    transition from any of the possible states.
\end{enumerate}
\begin{definition}[Language of an NFA]
  The \textit{language} of an NFA $A$ is denoted $L(A)$, and is defined by
  \[L(A)=\left\{ \omega|\hat{\delta}\left( q_0,\omega \right) \cap F\neq \emptyset \right\},\]
  the set of string $\omega$ that the start state may take that does not
  definitively lead to a nonaccepting state. The set of strings such that the
  possibilities contain at least one accepting state.
\end{definition}

\subsubsection{NFAs with $\epsilon$-Transitions}
We consider an automaton in which we allow a transition on $\epsilon$. The NFA
is hence allowed to make a transition without receiving an input symbol.
\begin{figure}[h!]
  \centering
  \includegraphics[width=0.5\textwidth]{nfa-epsilon-diag}
\end{figure}
\begin{definition}[$\epsilon$-NFA]
  An $\epsilon$-NFA $A=\left( Q,\Sigma,\delta,q_0,F \right) $, where:
  \begin{itemize}
    \item $Q$: A finite set of states
    \item $\Sigma$: A finite set of input symbols
    \item $\delta$: A transition function taking as argument a state and an
      input symbol (including $\epsilon$) and returns \textbf{a set of states}.
      $\delta: Q\times (\Sigma\cup \left\{ \epsilon \right\})\to \mathcal{P}(Q)$
    \item $q_0$: A start state, $q_0\in Q$
    \item $F$: A set of final or accepting states. $F\in \mathcal{P}\left( Q \right) $
  \end{itemize}
\end{definition}
To define an extended transition function for $\epsilon$-NFAs, we first define
the \textit{$\epsilon$-closure} of a state recursively:
\begin{enumerate}
  \item $q\in \epsilon$-closure$\left( q \right) $
  \item If $p\in \epsilon$-closure$\left( q \right)\wedge\delta\left( p,\epsilon \right) =r \implies r\in \epsilon$-closure$\left( q \right) $.
\end{enumerate}
The $\epsilon$-closure of a state $q$ is the set of all states reachable from
$q$ using only $\epsilon$ transitions.
We further define the $\epsilon$-closure of a set of states $S$:
\[\epsilon\text{-closure}\left( S \right) =\bigcup_{q\in S} \epsilon\text{-closure}\left( q \right).\]
We now define the \textbf{extended transition function}
$\hat{\delta}:Q\times \Sigma^{*}\to\mathcal{P}\left( Q \right) $, the set of
states that the $\epsilon$-NFA will be in if it starts in a given state $q$ and processes
the string $\omega$. We define by induction on the length of the input string:
\begin{enumerate}
  \item $\hat{\delta}\left( q,\epsilon \right) =\epsilon\text{-closure}\left(
    q \right)$.\ (by the definition of $\epsilon$-closure) \item Suppose $w$ a string of the form $xa$, where $x\in \Sigma^{*},a\in
    \Sigma$ (WLOG, $\epsilon$ not in $\Sigma$), and that $\hat{\delta}\left(
    q,x \right) =\left\{ p_1,p_2,\ldots,p_k \right\}$
    \[\hat{\delta}\left( q,w \right) =\bigcup_{i=1}^{k}\delta\left( p_i,a
    \right) =\epsilon\text{-closure}\left(   \left\{ r_1,r_2,\ldots,r_m \right\}\right)\]
    That is, we first compute $\hat{\delta}\left( q,x \right)$, then follow any
    transition from any of the possible states, then consider the possibility
    of any $\epsilon$ transitions from the possible states.
\end{enumerate}

\begin{remark}
  A $\epsilon$-NFA $E=(Q_E,\Sigma,\delta_E,q_0,F_E)$ accepts the same language
  as the DFA $D=\left( Q_D,\Sigma,\delta_D,q_D,F_D \right) $ where:
  \begin{itemize}
    \item $Q_D=\left\{ \epsilon\text{-closure}\left( S \right) | S\in
      \mathcal{P}(Q_E) \right\}$
    \item $q_D=\epsilon\text{-closure}\left( q_0 \right) $
    \item $F_D = \left\{ S| S\subseteq Q_D \wedge S\cap F_E\neq
      \emptyset\right\}$
    \item $\delta_D\left( S,a
  \right) =\epsilon\text{-closure}\left( \bigcup_{p\in S} \delta_E\left( p,a \right)  \right) $
  \end{itemize}
\end{remark}

\subsection{DFAs vs NFAs}
\subsubsection{Equivalence of DFAs and NFAs}
\begin{theorem}
  A language $L$ is accepted by some DFA if $L$ is accepted by some NFA\@.
\end{theorem}
\begin{proof}
Suppose an NFA $N=\left( Q_N,\Sigma,\delta_N,q_0,F_N \right) $. We consider the
DFA $D=\left( Q_D,\Sigma,\delta_D,\left\{ q_0 \right\} ,F_D \right) $, where
\begin{itemize}
  \item $Q_D =\mathcal{P}\left( Q_N \right) $. Intuitively, each state of the
    DFA represents a state of the entire NFA\@.
    \begin{remark}
      This conversion is known as \textit{subset construction}, involving
      constructing all subsets of the set of states of the NFA\@. Often, not all
      states are accessible from the start state. Inaccessible states can be
      discarded.
    \end{remark}
  \item $F_D = \left\{ S| S\subseteq Q_N \wedge S\cap F_N\neq
    \emptyset\right\}\subseteq \mathcal{P}(Q_N)$. That is, $F_D$ consists all
    sets of $N$'s states that include at least one accepting state of $N$.
  \item $\delta_D\left( S,a \right) =\bigcup_{p\in S} \delta_N\left( p,a
    \right) $, given $S\in \mathcal{P}\left( Q_N \right)$, and $ a\in \Sigma$.
    We consider all possible current states, and take the union of all possible
    states after processing input $a$
\end{itemize}
To prove that $\forall w\in \Sigma^{*}\left( \hat{\delta}_D\left( \left\{ q_0 \right\},\omega \right)=\hat{\delta}_N\left( q_0,w \right) \right) $, we induct on the length of $\omega$:
\begin{enumerate}
  \item \textit{Base case:} $\omega=\epsilon$
    \[\hat{\delta}_D\left( \left\{ q_0  \right\} ,\epsilon  \right)=\left\{ q_0
    \right\}=\hat{\delta}_N\left( q_0,\epsilon \right).\]
  \item \textit{Inductive step: }
    \begin{align*}
      \hat{\delta}_D\left( \left\{ q_0 \right\},\omega a\right)=&\delta_D\left(
      \hat{\delta}_D\left( \left\{ q_0 \right\},\omega \right) ,a \right)\\
        =&\bigcup_{q\in \hat{\delta}_D\left( \left\{ q_0 \right\} ,\omega\right)} \delta_N\left( q,a \right) \\
        =&\bigcup_{q\in \hat{\delta}_N\left(  q_0 ,\omega\right)} \delta_N\left( q,a \right) \\
        =&\hat{\delta}_N\left( q_0,wa \right).
    \end{align*}
\end{enumerate}
Since $\hat{\delta}_D\left( \left\{ q_0 \right\} ,\omega\right)\in F_D\iff\hat{\delta}_N\left( q_0,wa \right)\cap F_N \neq \emptyset$, we have $L(N)=L(D)$.\qed%
\end{proof}

\begin{theorem}
  A language $L$ is accepted by some DFA only if $L$ is accepted by some NFA\@.
\end{theorem}
\begin{proof}
  Intuitively, we take the transition diagram for the DFA as a transition diagram
  for an NFA, which has only one choice of transition at any state.\ \qed%
\end{proof}

Consider a language that can be described by some NFA and some DFA\@. The DFA in
practice has about as many states as the NFA, but often with more transitions.
In the worst case, however, the smallest DFA can have $2^{n}$ states while the
smallest NFA for the same language only has $n$ states.

\subsubsection{Languages of Automata}
\begin{definition}[Regular Language]
  If $L$ is $L(A)$ for some DFA $A$, then $L$ is a regular language.
\end{definition}
\begin{corollary}
  $L$ is a regular language $\iff$ $L$ is $L(N)$ for some NFA $N$.
\end{corollary}

\begin{theorem}
  Given a non-empty finite alphabet $\Sigma$, there exists a countable number of regular languages.
\end{theorem}
Since each regular language can be represented by an automaton, there exist
finite, or countably infinite regular languages. That is, there exist
languages that are not accepted by any automata and is hence non-regular.

\subsection{Minimisation of Automata}\label{minimisation}

\subsubsection{Prerequisites}
Suppose a regular language $L$. We define the equivalence relation for $u,w\in \Sigma$
\[u \equiv_L w\iff \left[ \forall x\in \Sigma, ux\in L \iff wx\in L \right].\]
We denote $equiv(w)$ the equivalence class of $w$.

\begin{definition}[Distinguishable States]
  We say that $(p,q)$ is distinguishable $\iff\exists w\in \Sigma^{*}$ such that
  \[\hat{\delta}\left( p,w \right)\in F\wedge \hat{\delta}\left( q,w \right)
  \notin F\qquad\vee\qquad\hat{\delta}\left( p,w \right)\notin F\wedge
\hat{\delta}\left( q,w \right) \in F.\] 
\end{definition}
\begin{remark}
  \[(p,q)\text{ is indistinguishable}\iff \left[ \forall w\in \Sigma^{*},
  \hat{\delta}\left( p,w \right)\notin F\iff \hat{\delta}\left( q,w \right) \in
F\right].\] 
\end{remark}

\subsubsection{Minimal DFA for $L$}
If there is a finite number of equivalence classes, $\exists$ a DFA $D=\left(
Q,\Sigma,\delta,q_0,F\right)$ where:
\begin{itemize}
  \item $Q=\left\{ equiv\left( \omega \right) | \omega\in \Sigma^{*} \right\}=\Sigma^{*}/\equiv_L$
  \item $\delta:(\Sigma^{*}/\equiv_L \times \Sigma)\to \Sigma^{*}/\equiv_L$ defined
    $\delta\left( equiv(\omega),a \right) =equiv(wa)$
  \item $q_0=equiv(\epsilon)$
  \item $F=\left\{ equiv(\omega) | w\in L \right\}$
\end{itemize}
which, by definition of the equivalence relation, accepts $L$. It is also the
minimal DFA for $L$.

\begin{lemma}
  $u \not\equiv_L \omega\implies \hat{\delta}'\left( q_0',u \right) \neq
  \hat{\delta}'\left( q_0',w \right) $.
\end{lemma}
\begin{proof}
  Suppose not. That is $u\not \equiv_L
  w\wedge\hat{\delta}'\left( q_0',u \right) = \hat{\delta}'\left( q_0',w
  \right)$. Then $\exists x\in \Sigma^{*}$ such that WLOG, $ux\in L\wedge
  wx\notin L$. But $\hat{\delta}'\left( q_0',u \right) = \hat{\delta}'\left(
  q_0',w \right)$ and $A'$ accepts either both or neither. Hence a
  contradiction.\qed%
\end{proof}

\begin{prop}
  $D$ is minimal.
\end{prop}
\begin{proof}
  Suppose a DFA $A'=\left( Q',\Sigma,\delta',q_0',F' \right) $ also accepting
  $L$. From the lemma, we have
  \[u \not\equiv_L \omega\implies \hat{\delta}'\left( q_0',u \right) \neq
  \hat{\delta}'\left( q_0',w \right).\]
  Let $u\equiv_{A'}w\iff\hat{\delta}'\left( q_0',u \right) =
  \hat{\delta}'\left( q_0',w \right)$. Then $u\equiv_L \implies u\equiv_{A'}
  w$. Hence $\equiv_{A'}$ divides the equivalence classes of $\equiv_{L}$ into
  finer equivalence classes, and $\equiv_L$ is minimal.\qed%
\end{proof}

\subsubsection{Minimal DFA Construction Algorithm}
Suppose we are given a DFA $A=\left( Q,\Sigma,\delta,q_0,F \right) $.
We determine all distinguishable pairs with the following recursive algorithm:
\begin{enumerate}
  \item \textit{Base case:} Each pair $(p,q)$ such that $p\in F$ and $q\notin
    F$ (or vice versa) is distinguishable.
  \item \textit{Inductive step:} If $\exists a\in \Sigma\left[  \delta\left(
    p,a \right) \text{ and }\delta\left( q,a \right) \text{ are
  distinguishable} \right]$, then $(p,q)$
    is distinguishable.
  \item Steps 1 and 2 are repeated until no more pairs can be added.
  \item The remaining pairs are non-distinguishable states.
\end{enumerate}
To compute the distinguishable states, we utilise the table-filling algorithm:
\begin{figure}[h!]
  \centering
  \includegraphics[width=0.3\textwidth]{table-filling}
\end{figure}
We then form the new DFA as follows:
\begin{enumerate}
  \item Remove all unreachable states.
  \item Find all nondistinguishable pairs of states using the algorithm. They
    are equivalent. This yields an equivalence relation
  \item Form the new DFA with:
    \begin{itemize}
      \item $Q_{new}$ as the set of the equivalence classes.
      \item The same $\Sigma$
      \item The transitions from each equivalence class is based on the
        corresponding transition in the original DFA\@:
        \[\delta\left( p,a \right) =q \implies\delta_{new}\left( E_p,a \right)=E_q.\]
        where $E_p, E_q$ are the equivalence classes of $p,q$ respectively.
      \item $q_{0,new}=equiv(q_0)$
      \item $F_{new}=\left\{ equiv(\omega) | w\in F \right\}$
    \end{itemize}
\end{enumerate}
% TODO: Add proof of correctness of table-filling algorithm
% TODO: Add proof of correctness of new DFA

\subsection{Parallel Simulation of n-DFAs}
Suppose two DFAs $A=\left(Q,\Sigma,\delta,q_0,F\right)$ and $A'=\left( Q',\Sigma,
\delta',q_0',F' \right) $.
We may construct another DFA $A''=\left( Q\times
Q',\Sigma,\delta'',\left(q_0,q_0'\right),F'' \right) $, where
\[\delta''\left( \left( q,q' \right) ,a \right) =\left( \delta\left( q,a
\right) ,\delta'\left( q',a \right)  \right).\] 
and $F''$ depends on need.
\begin{example}[Definitions of $F''$]\ 
   \begin{itemize}
    \item $F''=F\times F'$ finds the intersections of the two automata.
    \item $F''=F\times Q'\cup Q\times F'$ finds the union of the two automata.
  \end{itemize}
\end{example}
This simulates the two DFAs in parallel.
\begin{remark}
  By repeated application, we may simulate any boolean combination of
  automata.
\end{remark}


\subsection{Regular Expressions}
Through regular expressions, another type of language-defining notation that
algebraically describes a language, we are able to declaratively express the
strings we want to accept. Each regular expression denotes a (regular) language.
\subsubsection{Operators}
Given a regular expression $E$, we denote $L(E)$ for the language that $E$
denotes. We then define regular expressions inductively (with various
operators), as follows:
\begin{enumerate}
  \item \textit{Basis Step:}
    \begin{itemize}
      \item Constant $\epsilon$ is a regular expression, with
        $L(\epsilon)=\left\{ \epsilon \right\}$.
      \item Constant $\emptyset$ is a regular expression, with $L(\emptyset)=\emptyset$.
      \item If $a\in \Sigma$, $a$ is a regular expression, with $L(a)=\left\{ a
        \right\}$.
    \end{itemize}
  \item \textit{Inductive Step:} Given regular expressions $E$ and $F$,
    \begin{itemize}
      \item $E+F$ is a regular expression denoting union, $L(E+F)=L(E)\cup L(F)$.
      \item $EF$ is a regular expression denoting concatenation, $L(EF)=L(E)L(F)$.
      \item  $E^{*}$ is a regular expression denoting its (Kleene's) closure,
        $L(E^{*})={\left( L\left( E \right)  \right)}^{*}$
      \item $E^{+}$ is a regular expression denoting its closure, without $\epsilon$, $L(E^{+})=EE^{*}$
      \item $(E)$ is a regular expression denoting the same language, $L\left( \left( E \right)  \right) =L(E)$
    \end{itemize}
\end{enumerate}
We further define the order of precedence of the operators:
 \begin{enumerate}
  \item The parantheses operator is of highest precedence, to specify the
    desired grouping of regular expressions.
  \item The star operator is of highest precedence, applied to the smallest
    sequence of symbols to its left that is a well-formed regular expression.
  \item The concatenation operator groups all expressions that are
    \textit{juxtaposed} (adjacent without intervening operators). Concatenation is associative.
  \item The union binary operator groups its operands. Union is associative.
\end{enumerate}

\subsubsection{Converting Automaton to Regular Expression}
\begin{theorem}  
  If $L = L(A)$ for some DFA $A$, then there is a regular expression $R$ such
  that $L=L(R)$. That is, every regular language can be defined by a regular
  expression.
\end{theorem}
\begin{proof}
  Suppose a DFA $A=\left( Q,\Sigma,\delta,q_{start},F \right)$. Without loss of generality, suppose
  $Q=\left\{ 1,2,\ldots,n \right\}$ for some $n$ and $q_{start}=1$.
  Let $R^{k}_{i,j}$ denote the regular expression for the set of strings which
  can be formed by going from state $i$ to state $j$ only using intermediate
  states $\le k$. We have
  \[L(A)=L\left(\sum_{j\in F} R_{1,j}^{n}\right)\]
  We compute by induction:
  \begin{enumerate}
    \item \textit{Base step:} $R_{i,j}^{0}$. That is, the path must have no
      intermediate states at all. We only have an arc from (state) $i$ to $j$,
      or a path of length 0 consisting only $i$.\\
      \[i\neq j \implies R_{i,j}^{0}=a_1+a_2+\cdots+a_m, \qquad\forall a_r\in \Sigma, \delta\left( i,a_r \right) =j.\] 
      \[i=j \implies R_{i,i}^{0}=\epsilon+a_1+a_2+\cdots+a_m, \qquad\forall a_r\in \Sigma, \delta\left( i,a_r \right) =i.\] 
    \item \textit{Inductive step:} Suppose there is a path from state $i$ to
      $j$ going through states no higher than $k$.
      \[R_{i,j}^{k+1}=R_{i,j}^{k}+R_{i,k+1}^{k}{\left(R_{k+1,k+1}^{k}\right)}^{*}R_{k+1,j}^{k}.\] 
      We then have the case where the path does not go through state $k+1$ at
      all, and the case where it goes through $k+1$ at least once.
  \end{enumerate}
  \begin{figure}[h!]
    \centering
    \includegraphics[width=0.5\textwidth]{automata-regex}
  \end{figure}
  We may then obtain
  $L(A)=L\left(\sum_{j\in F} R_{1,j}^{n}\right)$.\qed%
\end{proof}
\begin{remark}
  The construction of a regular expression from automata could be applied to
  NFA or $\epsilon$-NFA, but we construct about $n^{3}$ expressions, with
  each's length growing by a factor of 4 each inductive step, resulting in a
  expression on the order of $4^{n}$ symbols for an $n$-state automaton.
\end{remark}
% TODO: Add Converting DFA's to Regular Expressions by Eliminating States TB pg. 98

\subsubsection{Converting Regular Expression to Automaton}
\begin{theorem}
  Every language defined by a regular expression can be defined by an automaton.
\end{theorem}
We may construct an automaton through structural induction on the expression,
constructing automata for the basis expressions and combining them into
larger automata that fits the various operators.

\begin{proof}
  Suppose $L=L(R)$for a regular expression $R$. We construct a $e$-NFA $E$ with:
  \begin{enumerate}
    \item Exactly one accepting state.
    \item No transition into the starting state.
    \item No transition out of the accepting state.
    \item Starting state $\neq$ Accepting state.
  \end{enumerate}
  For the bases, we create automata that accepts the languages 
  \begin{itemize}
    \item (a) $L(\epsilon)$,
    \item (b) $L(\emptyset)$,
    \item (c) $L(a)$
  \end{itemize}
  We then assume that we have the automaton for the immediate subexpressions of
  a given regular expression. We then have the four cases: 
  \begin{itemize}
    \item (a) $R+S$,
    \item (b) $RS$,
    \item (c) $R^{*}$
    \item $R^{+}$ which is constructed from the automaton of $R^{*}$, but
      without the $\epsilon$ transition from the starting to the accepting
      state
    \item $(R)$ where the automaton does not change.
  \end{itemize}
  \begin{figure}[h!]
    \centering
    \raisebox{-0.5\height}{\includegraphics[width=0.2\textwidth]{regex-automata}}
    \hspace*{.2in}
    \raisebox{-0.5\height}{\includegraphics[width=0.35\textwidth]{regex-automata-2}}
  \end{figure}
  The constructed automata satisfy the four conditions, providing the language
  of the regular expression.\qed%
\end{proof}

\subsubsection{Algebra of Regular Expressions}
We find various useful results for simplifying regular expressions and to keep
the size of expressions manageable. We consider the algebra of regular
expressions, analysing properties such as associativity, commutativity and
distributivity, as well as relevant identities and annihilators. For regular
expressions $L, M, N$:
\paragraph{Associativity, Commutativity and Distributivity}
\begin{itemize}
  \item $L+M=M+L$: Commutative law for union.
  \item  $(L+M)+N=L+(M+N)$: Associative law for union. With the commutative
    law, we may take union in any order.
  \item $(LM)N=L(MN)$: Associative law for concatenation. Concatenation is non-commutative.
  \item $L(M+N)=LM+LN$: Left-distributive law of concatenation over union.
  \item $(M+N)L=ML+NL$: Right-distributive law of concatenation over union.
\end{itemize}
\paragraph{Identities, Annihilators and Idempotence}
\begin{itemize}
  \item $\emptyset+L=L+\emptyset=L$: $\emptyset $ identity for union.
  \item $\epsilon L = L \epsilon=L$: $\epsilon$ identity for concatenation.
  \item $\emptyset L=L\emptyset =\emptyset $: $\emptyset $ annihilator for concatenation.
  \item $L + L = L$: Idempotence law for union.
\end{itemize}
\paragraph{Closures}
\begin{itemize}
  \item ${\left( L^{*} \right)} ^{*}=L^{*}$: Closing an expression that is
    already closed does not change.
  \item $\emptyset ^{*}=\epsilon$
  \item $\epsilon^{*}=\epsilon$ 
  \item $L^{+}=LL^{*}=L^{*}L$ 
  \item $L^{*}=L^{+}+\epsilon$
  \item ${(L+M)}^{*}={(L^{*}M^{*})}^{*}$
\end{itemize}

\subsection{Regular Languages}
To show that a given language is nonregular, we may use the pumping lemma.
Intuitively, we can always find a nonempty string $y$ not too far from the
beginning that can be repeated any number of times, or deleted while keeping
the resulting string in $L$.
\begin{lemma}
  Let $L$ be a regular language. Then $\exists n=f(L)$ such that $\forall\omega\in L$,
  $\left| \omega \right| \ge n$, $\exists x,y,z\in \Sigma^{*},\omega=xyz$ such that:
  \begin{itemize}
    \item $\left| xy \right| \le n$. That is, string to be pumped is not too
      far from the start of the string.
    \item $y\neq\epsilon$. The string to be ``pumped'' cannot be empty.
    \item $\forall i\ge 0$, $xy^{k}z\in L$
  \end{itemize}
\end{lemma}
The contrapositive of the lemma is used to prove that a language is
nonregular. That is, $\forall n\in \mathbb{N}$, $\exists w\in L$, such that
$\left| \omega \right| \ge n$, $\forall x,y,z\in \Sigma^{*}$, where $\omega=xyz$,
$\left| xy \right| \le n$, $y\neq\epsilon$, we have
\[\exists i\in \mathbb{N}\left( xy^{i}z\notin L \right)\]
\begin{proof}
  Suppose $L$ is regular. Then $\exists$ DFA $A=\left( Q, \Sigma, \delta, q_0,F
  \right)$ such that $L=L(A)$. Suppose $\left| Q \right|=n$. Consider a string
  $\omega\in \Sigma^{*},\left| w \right| \ge n$. Then $w=a_1a_2\cdots a_m$ for
  some $m\ge n, a_i\in \Sigma$. For $i=0,1,\ldots,n, p_i:=\hat{\delta}\left(
  q_0,a_1a_2\cdots a_i \right)$. $p_i$ is the state of $A$ after reading the
  first $i$ symbols of $\omega$.\\
  Since there are $n$ different states and $n+1$ different $p_i$'s, by the
  pigeonhole principle, $\exists 0\le i<j<n$ such that $p_i=p_j$. Let:
  \begin{itemize}
    \item $x=a_1a_2\cdots a_i$ 
    \item $y=a_{i+1}a_{i+2}\cdots a_j$
    \item $z=a_{j+1}a_{j+2}\cdots a_m$
  \end{itemize}
  By processing $x$, $A$ takes up state $p_i$, after which $y$ takes it from
  $p_i$ back to $p_i$, and $z$ is the balance of $\omega$. $x=\epsilon\iff
  i=0$, and $z=\epsilon\iff j=n=m$. However, $y$ cannot be empty, since $i<j$.
  \begin{figure}[h!]
    \centering
    \includegraphics[width=0.5\textwidth]{rl-pumping-lemma}
  \end{figure}
  Then consider for any $k\in \mathbb{N}$, the string $xy^{k}z$:
  \begin{itemize}
    \item If $k=0$, the automaton goes from $q_0=p_0$ to $p_i=p_j$ on input
      $x$, and then from $p_i=p_j$ to the accepting state on input $z$. $A$
      accepts $xz$.
    \item If  $k>0$, the automaton goes from $q_0$ to $p_i$ on input $x$, and
      then circles $p_i$ to $p_i$ $k$ times on input $y^{k}$, and then to the
      accepting state on input $z$. $A$ accepts $xy^{k}z$.
  \end{itemize}\vspace{-7mm}\qed%
\end{proof}

By using the pumping lemma, we can construct a proof of the following structure
to show a specified language is non-regular:
\begin{prop}
  $L=\left\{ a^{m}b^{m}|m\ge 1 \right\}$ is nonregular.
\end{prop}
\begin{proof}
  Suppose $L$ is regular. Let $n$ be as in the Pumping Lemma. Let
  $\omega=a^{n}b^{n}$. Let $\omega=xyz$ as in the pumping lemma. Then $y$
  consists of only $a$'s, and $\left| y \right| >0$, $y={\left\{ a
  \right\}}^{+}$. By the pumping lemma, $xy^{2}z\in L$. But $xy^{2}z$ contains
  more $a$'s than $b$'s, and $xy^{2}z\notin L$. Hence $L$ is nonregular.
\end{proof}
\begin{prop}
  $L=\left\{ a^{p}|p\text{ is prime} \right\}$ is nonregular.
\end{prop}
\begin{proof}
  Suppose $L$ is regular. Let $n$ be as in the Pumping Lemma. Let
  $\omega=a^{p}$, where $p$ is prime and $p>n$. Let $\omega=xyz$ as in the
  pumping lemma. By the pumping lemma, $\forall k\in \mathbb{N}, xy^{k}z\in L$. Consider
  $k=p+1$. Then $xy^{k}z=a^{r}$ where 
  \[r=\left| x \right| +\left| z \right| +\left( p+1 \right) \left| y \right| =p+p\left| y \right|  \]
  and $r$ is non-prime,  and $xy^{p+1}z\notin L$. Hence $L$ is nonregular.
\end{proof}
\begin{prop}
  $L= \left\{a^i b^j | i < j\right\}$ is nonregular.
\end{prop}

\subsubsection{Closure of Regular Languages}
We have certain operations such that, when applied to the languages, the
resulting related language is also regular.

\begin{definition}[Homomorphism]
  Suppose two alphabets $\Sigma$ and $\Gamma$. Consider a mapping
  $h:\Sigma\to\Gamma^{*}$. We further define $h:\Sigma^{*}\to\Gamma^{*}$ by
  principle of induction:
  \begin{enumerate}
    \item \textit{Base case:}
      $h(\epsilon)=\epsilon$
    \item \textit{Inductive Step:}
      $h(a\omega)=h(a)\cdot h(\omega), \forall a\in \Sigma,\omega\in \Sigma^{*}$ 
  \end{enumerate}
  $h$ is called a \textit{homomorphism}. Intuitively, a homomorphism maps each
  symbol to another symbol in every string of the language.
\end{definition}

Given regular languages $L, L_1,L_2$:
\begin{itemize}
  \item $L_1\cup L_2$ is regular.
  \item $L_1\cdot L_2$ is regular.
  \item $\bar{L}=\Sigma^{*}-L$ is regular.
  \item $L_1\cap L_2$ is regular.
  \item $L_1-L_2$ is regular.
  \item $L^{R}$ is regular.
  \item $h(L)$ is regular, for some homomorphism $h$.
\end{itemize}
We may prove the results by constructing a regular expression through principle
of induction on its length, from the base cases of a regular expression
$\emptyset,\epsilon,a\in \Sigma$ and the defined operations.
\begin{remark}
  Supposing $L$ is a regular language over an alphabet $\Sigma_1$. Then $L$ is also a regular language
  over an alphabet $\Sigma_2\supseteq\Sigma_1$, by considering the same DFA over $\Sigma_2$.
\end{remark}

\subsubsection{Decision Problems of Regular Languages}
We may examine some properties of regular languages through conversion between
the equivalent forms of representing regular languages.\\

\begin{figure}[h!]
  \centering
  \includegraphics[width=0.3\textwidth]{conversions}
\end{figure}

We check the \textbf{emptiness of regular languages} by considering the automaton, and
checking whether any accepting states are reachable from the starting state. That is,
\[L=\emptyset \iff \forall q\in F, q\text{ is unreachable}.\] 
We compute this using a recursive algorithm:
\begin{itemize}
  \item \textit{Base case:} The start state is reachable from the start state.
  \item \textit{Inductive step:} If $p$ is reachable from the start state, all
    states $q$ such that $\exists a\in \Sigma, \delta(p,a)=q$ are reachable
    from the start state.
\end{itemize}

We check for \textbf{membership} by converting whichever representation available into a
DFA, and then simulating the DFA\@. Alternatively, if given another automaton, we
may attempt to simulate the NFA or $\epsilon$-NFA itself, keeping track of all
possible states.\\

\textbf{Two automata $A, A'$ are equivalent} $\iff$ the minimisation of $A$ is the same
as the minimisation of $A'$, except for renaming of the states. (Minimisation
of DFA, Subsection~\ref{minimisation})

\subsection{Linear Grammars}
\begin{definition}[Right Linear Grammar]
  A \textbf{Right Linear Grammar} is a CFG where all productions are of the form:
  \begin{align*}
  A\to\omega\beta\qquad &\omega\in T^{*},\quad \beta\in V\\
  A\to \omega\qquad&\omega\in T^{*}
  \end{align*}
  Similarly, a \textbf{Left Linear Grammar} is one in which all productions are of the form:
  \begin{align*}
  A\to\beta\omega\qquad &\omega\in T^{*},\quad \beta\in V\\
  A\to \omega\qquad&\omega\in T^{*}
  \end{align*}
\end{definition}
\begin{theorem}
  Every regular language can be generated by a right-linear grmmar.
\end{theorem}
\begin{proof}
  Suppose a regular language $L$. There exists a DFA $A=\left( Q,\Sigma,\delta,q_0,F \right) $ such
  that $L(A)=L$. WLOG, assume that $Q\cap \Sigma=\emptyset$. Consider the
  right-linear grammar $G=\left( Q,\Sigma,P,q_0 \right) $, where:
  \begin{enumerate}
    \item $\forall q,p\in Q,a\in \Sigma$
      \[\delta\left( q,a \right) =p\iff\exists\text{ a production }(q\to ap)\in P.\] 
    \item $\forall q\in F$
      \[\exists \text{ a production }(q\to\epsilon)\in P.\]
  \end{enumerate}
  We induct on the length $\left| \omega \right| $ to show that
  $\hat{\delta}\left( q_0,\omega \right)=p\iff q_0 \Rightarrow _G^{*}\omega p$.
  \begin{enumerate}
    \item \textit{Basis Step}: For $\left| \omega \right| =0$, $\omega=\epsilon$ and 
      \[\hat{\delta}\left( q_0,\omega \right)=p\iff p=q_0\iff q_0 \Rightarrow _G^{*}\omega p\]
    \item \textit{Inductive Step}: Consider $\omega=ua$, for some $a\in \Sigma$
      and $u\in \Sigma_{\le k}$. We have:
      \begin{align*}
        \hat{\delta}\left( q_0,ua \right) =p'\iff&\exists p\in Q,\quad
        \hat{\delta}\left( q_0,u \right) =p\wedge\delta\left( p,a \right) =p'\\
        \iff&\exists p\in Q,\quad q_0 \Rightarrow _G^{*}up\wedge p\to ap'\\
        \iff&q_0 \Rightarrow _G^{*}uap'=\omega p'
      \end{align*}\vspace{-4mm}\qed%
  \end{enumerate}
\end{proof}
\vspace{-4mm}
\begin{theorem}
  A language generated by a right-linear grammar is regular.
\end{theorem}
\begin{proof}
  Suppose a right-linear grammar $G=\left( V,\Sigma,P,S \right)$. WLOG, suppose
  $V\cap \Sigma=\emptyset$ and each productions are of the form $A\to bC$ or $A\to\epsilon$ where
  $A,C\in V$ and $b\in \Sigma\cup \left\{ \epsilon \right\} $. Consider the
  NFA $A=\left( V,\Sigma,\delta,S,F \right) $, where:
  \[\left( A\to aB \right)\in P\iff B\in \delta\left( A,a \right).\]
  and $F =\left\{ A\in V|(A\to\epsilon)\in P \right\} $
  We induct on the length $\left| \omega \right| $ to show that $A \Rightarrow _G^{*}\omega B\iff B\in \hat{\delta}\left( A,\omega \right) $.
  \begin{enumerate}
    \item \textit{Basis Step}: For $\left| \omega \right| =0$, $\omega=\epsilon$ and
      \[A \Rightarrow _G^{*}B\iff A=B \iff B\in \hat{\delta}\left( A,\epsilon \right).\]
    \item Consider $\omega=ua$. Then
      \begin{align*}
        A \Rightarrow _G^{*} \omega B=uaB\iff&\exists C\in V,\quad A\Rightarrow _G^{*}uC\wedge C\Rightarrow _G aB\\
        \iff&\exists C\in V,\quad\hat{\delta}\left( A,u \right) =C\wedge \delta\left( C,a \right) =B\\
        \iff&\hat{\delta}\left( A,ua \right)=\hat{\delta}\left( A,\omega \right)  =B
      \end{align*}\vspace{-4mm}\qed%
  \end{enumerate} 
\end{proof}


